\documentclass[border=8pt]{standalone}
\usepackage[T1]{fontenc}
\usepackage[utf8]{inputenc}
\usepackage[spanish]{babel}
\usepackage{lmodern}
\usepackage{tikz}
\usetikzlibrary{arrows.meta,positioning,babel}
\definecolor{azul}{HTML}{1F4E79}
\definecolor{naranja}{HTML}{C9601F}
\definecolor{verde}{HTML}{2E7D32}
\definecolor{rojo}{HTML}{B23A34}
\begin{document}
\begin{tikzpicture}[font=\small,>=Latex,
 pin/.style={draw,fill=white,minimum width=1.5cm,minimum height=0.55cm,font=\footnotesize\bfseries},
 pc/.style={draw=azul,thick,rounded corners,fill=naranja!8,align=center,text width=3.0cm,minimum height=1.1cm}]
% --- Tarjeta ESP32-S3 ---
\draw[azul,very thick,rounded corners,fill=azul!6] (0,0) rectangle (4.2,5.2);
\node[azul,font=\bfseries,align=center] at (2.1,4.7) {ESP32-S3-DevKitC-1};
\node[pin] (e3v3) at (3.4,3.6) {3V3};
\node[pin] (egnd) at (3.4,2.7) {GND};
\node[pin] (e8)   at (3.4,1.8) {GPIO8};
\node[pin] (e9)   at (3.4,0.9) {GPIO9};
\node[font=\footnotesize,align=center,text width=1.9cm] at (1.2,2.2) {SDA = 8\\SCL = 9\\(libres junto al display de la P4)};
% --- Modulo ZS-042 ---
\draw[verde,very thick,rounded corners,fill=verde!6] (11.0,0) rectangle (15.6,5.2);
\node[verde,font=\bfseries,align=center] at (13.3,4.7) {Módulo DS3231 (ZS-042)};
\node[pin] (mvcc) at (11.8,3.6) {VCC};
\node[pin] (mgnd) at (11.8,2.7) {GND};
\node[pin] (msda) at (11.8,1.8) {SDA};
\node[pin] (mscl) at (11.8,0.9) {SCL};
\node[font=\footnotesize,align=center,text width=2.6cm] at (14.0,2.3) {DS3231 en \textbf{0x68}\\EEPROM AT24C32 en \textbf{0x57}\\pull-ups de 4.7\,k\\a VCC ya incluidas};
\draw[fill=gray!25] (14.0,0.75) circle (0.45);
\node[font=\scriptsize,align=center] at (14.0,0.75) {pila\\2032};
% --- Cables ---
\draw[rojo,ultra thick] (e3v3.east) -- node[above,font=\footnotesize]{rojo: 3.3\,V (NO 5\,V)} (mvcc.west);
\draw[black,ultra thick] (egnd.east) -- node[above,font=\footnotesize]{negro: GND común} (mgnd.west);
\draw[azul,ultra thick] (e8.east) -- node[above,font=\footnotesize]{SDA: datos (bidireccional)} (msda.west);
\draw[naranja,ultra thick] (e9.east) -- node[above,font=\footnotesize]{SCL: reloj (lo genera el ESP32-S3)} (mscl.west);
% --- Laptop ---
\node[pc,left=1.4cm of e3v3.west,anchor=east,yshift=-1.4cm,xshift=-2.6cm] (pc) {Laptop\\Monitor PlatformIO\\115200};
\draw[<->,thick] (pc.east) -- node[above,font=\footnotesize]{USB} (0,2.2);
% --- Osciloscopio (opcional) ---
\node[draw,dashed,rounded corners,align=left,font=\footnotesize,text width=4.6cm] (osc) at (7.6,-1.6)
 {\textbf{Osciloscopio / analizador (Paso 7):}\\CH1 en SDA, CH2 en SCL, tierra de la sonda en GND. Disparo: flanco de bajada de SDA.};
\draw[dashed,->] (osc.north) -- (7.6,1.8);
\draw[dashed,->] ([xshift=0.8cm]osc.north) -- (8.4,0.9);
% --- Notas ---
\node[anchor=north west,text width=19cm,align=left,font=\footnotesize] at (-4.4,-2.9)
 {Cuatro cables Dupont hembra--hembra; el módulo trae su conector de 6 pines (32K y SQW no se usan hoy).
  Alimenta el módulo con \textbf{3V3}: sus pull-ups van a VCC, y con 5\,V llevarían SDA/SCL a 5\,V, más de lo que tolera el ESP32-S3 (3.6\,V).
  Con pila \textbf{CR2032} (no recargable) a 3.3\,V el circuito de carga del módulo casi no conduce; con 5\,V puede dañarla.};
\node[anchor=north west,text width=19cm,align=left,text=naranja,font=\footnotesize] at (-4.4,-4.1)
 {Conecta y desconecta cables con el USB desconectado. Si \texttt{scan} no encuentra nada: revisa VCC, GND y que SDA vaya a 8 y SCL a 9 (no cruzados).
  Pruebas del TC07 que se hacen a propósito: quitar el cable SDA (NACK), cambiar la dirección con \texttt{dir 0x69} y quitar VCC con y sin pila (bit OSF).};
\end{tikzpicture}
\end{document}
